\chapter{Lattices and sphere packings}\label{geometry:ch:lattices}

\begin{center}\itshape
\dots the packing will be the tightest possible, so that in no other arrangement could more pellets be stuffed into the same container.
\upshape(J.~Kepler, \emph{The Six-Cornered Snowflake}, 1611, as quoted by T.~C.~Hales)
\end{center}

Put equal coins on a table without overlapping them. A square grid leaves
visible gaps. Shift alternate rows into the gaps, and six coins can touch
one coin instead of four. The covered fraction rises from $\pi/4$ to
$\pi/\sqrt{12}$. Why can no other lattice arrangement improve it?

Stacking balls gives a harder version of the question. Harriot considered
it after Raleigh asked him how to count cannonballs in a pile. Kepler's
1611 essay on the snowflake proposed the familiar close packing as the
best arrangement. Gauss proved its optimality among lattice packings in
1831. Hales's proof for all three-dimensional packings, published in 2005,
required computer verification. In dimension eight, Viazovska proved in
2016 that the lattice $E_8$ is optimal among all packings; her work earned
a Fields Medal in 2022. These optimality theorems are quoted here.

Our proofs concern the arithmetic inside a lattice. A short nonzero
vector gives a nearby centre, so its length controls how large the balls
can be. A determinant gives the volume per centre. Minkowski's theorem
relates the two and also proves that every prime congruent to $1$ modulo
$4$ is a sum of two squares. In the plane, reducing a basis sharpens the
bound enough to settle the lattice packing problem exactly.

The appendix uses determinants and Euclidean volume, finitely generated
abelian groups from Volume~3, Chapter~7, and quadratic residues from
Volume~2, Chapter~6. Dimension $n$ is a positive integer whenever a shortest
nonzero lattice vector is discussed.

\begin{figure}[ht]
\centering
\fig{geometry-ch09-packings}
\caption{Disks of diameter $1$ centred at the points of the square lattice and of the hexagonal lattice. In rust, a fundamental parallelogram: each
contains exactly one disk's worth of area, $\frac\pi4$, and has area $1$ for the square lattice and $\frac{\sqrt3}2$ for the hexagonal one.}
\label{geometry:fig:packings}
\end{figure}

\begin{think}
The two fundamental parallelograms in Figure~\ref{geometry:fig:packings}
have areas $1$ and $\sqrt3/2$. Each accounts for one disc of area $\pi/4$.
Compute the two densities. Which part of that calculation depends on the
choice of basis, and which part should belong to the lattice itself?
\end{think}

\section{Lattices}

\begin{definition}
A full lattice in $\R^n$ is a subgroup
$\Lambda=\Z b_1+\cdots+\Z b_n$ generated by a real vector-space basis.
With basis matrix $B=(b_1\ \cdots\ b_n)$, it is $B\Z^n$.
Its Gram matrix is $G=B^{\mathsf T}B=(b_i\cdot b_j)$.
Its covolume is $\det\Lambda=|\det B|=\sqrt{\det G}$.
\end{definition}

The notation $\det\Lambda$ is positive, even if the determinant of a
chosen basis matrix is negative. For example, $(1,0),(0,1)$ and
$(1,0),(1,1)$ generate the same lattice $\Z^2$. The second parallelogram
is sheared, but still has area one.

\begin{proposition}\label{geometry:prop:latticebases}
Two real basis matrices $B,B'$ generate the same lattice if and only if
$B'=BU$ for an integer matrix $U$ with determinant $1$ or $-1$.
Consequently the covolume does not depend on the lattice basis.
\end{proposition}
\begin{proof}
If the lattices agree, each column of $B'$ is an integer combination of
columns of $B$, so $B'=BU$ for an integer matrix $U$.
The reverse inclusion similarly gives $B=B'V$ with $V$ integral.
Since $B$ is invertible, $UV=I$ and
$\det U\det V=1$. Both factors are integers, so $\det U=\pm1$.
Conversely, such a $U$ has integer inverse by the adjugate formula.
It maps $\Z^n$ onto itself, giving $B'\Z^n=B\Z^n$.
Finally $|\det B'|=|\det B||\det U|=|\det B|$.
\end{proof}

The half-open parallelepiped
\[
 \Pi=\{t_1b_1+\cdots+t_nb_n:0\le t_i<1\}
\]
has volume $\det\Lambda$. Every point of $\R^n$ lies in exactly one
translate $\lambda+\Pi$: split each basis coordinate into its integer
part and its fractional part. The half-open convention prevents overlap
on the boundary. Thus covolume is the volume assigned to one lattice point.

A bounded set contains only finitely many lattice points, since $B^{-1}$
takes it to a bounded set containing only finitely many integer vectors.
The lattice is therefore discrete. In dimension zero it is just $\{0\}$,
with covolume one; the strict volume hypotheses below cannot hold there.
For $n\ge1$ there is a shortest nonzero vector: choose any nonzero vector first and minimise among the finitely many
vectors no longer than it. Denote the resulting positive length by
$\lambda_1(\Lambda)$.

\section{Minkowski's theorem}

If a bounded region has more volume than a fundamental cell, moving its
pieces into one cell must make two pieces overlap. The resulting two
points differ by a lattice vector.

\begin{lemma}[Blichfeldt]\label{geometry:lem:blichfeldt}
If a bounded measurable set $S\subset\R^n$ has volume greater than
$\det\Lambda$, then $S$ contains distinct points $x,y$ with
$x-y\in\Lambda$.
\end{lemma}
\begin{proof}
Cut $S$ into the disjoint pieces $S_\lambda=S\cap(\lambda+\Pi)$.
Only finitely many are nonempty because both $S$ and $\Pi$ are bounded.
Their translates $S_\lambda-\lambda$ all lie in $\Pi$, but their
volumes sum to more than $\operatorname{vol}\Pi$.
They cannot be pairwise disjoint. Thus for some $\lambda\ne\mu$,
$x-\lambda=y-\mu$ with $x\in S_\lambda$, $y\in S_\mu$.
It follows that $x-y=\lambda-\mu\ne0$ is in $\Lambda$.
\end{proof}

A set is symmetric about $0$ if $x\in C$ implies $-x\in C$.
For a convex symmetric set, the difference of two points in half the set
belongs to the original set. This turns Blichfeldt's difference into a
point whose location we control.

The volume in this statement is well-defined. A convex set with empty
interior lies in a proper affine subspace: otherwise it contains an
$n$-dimensional simplex. It therefore has Lebesgue measure zero.
For a bounded convex set with interior, its closure $K$ is compact.
Translate an interior point to $0$. Choose a ball $B(0,r)$ already inside the original set $C$.
For $x\in K$, approximate $x$ by $x_j\in C$. If
$|y-(1-\varepsilon)x|<\varepsilon r/2$, then for sufficiently large $j$
we can write $y=(1-\varepsilon)x_j+\varepsilon b$ with $|b|<r$.
Convexity gives $y\in C$, proving
$(1-\varepsilon)K\subset\operatorname{int}C$.
Thus $K\setminus\operatorname{int}C$ has measure at most
$|K|-(1-\varepsilon)^n|K|$, which tends to zero.
Since $\operatorname{int}C\subset C\subset K$, completeness of
Lebesgue measure makes $C$ measurable, regardless of which boundary
points it includes.

\begin{theorem}[Minkowski]\label{geometry:thm:minkowski}
Let $C\subset\R^n$ be bounded, convex and symmetric about $0$.
If $\operatorname{vol}(C)>2^n\det\Lambda$, then $C$ contains a
nonzero point of $\Lambda$.
\end{theorem}
\begin{proof}
Convex sets are Lebesgue measurable; a lower-dimensional one has volume
zero, so the volume hypothesis in particular excludes that case.
The set $C/2$ has volume greater than $\det\Lambda$.
By Blichfeldt, it contains $x\ne y$ with $x-y\in\Lambda$.
The points $2x$ and $-2y$ belong to $C$. Their midpoint $x-y$ belongs
to $C$ by convexity and is a nonzero lattice point.
\end{proof}

The strict inequality cannot be replaced by a non-strict one for every
bounded convex set. The open cube $(-1,1)^n$ has volume $2^n$ but no
nonzero integer point. Compactness changes the boundary case, as
Exercise~\ref{geometry:ex:v9-closed-minkowski} shows.

\begin{theorem}[Fermat]\label{geometry:thm:twosquares}
Every prime $p\equiv1\pmod4$ can be written as a sum of two integer squares.
\end{theorem}
\begin{proof}
Quadratic residues from Volume~2, Chapter~6, give an integer $m$ with
$m^2\equiv-1\pmod p$. Consider
\[
 \Lambda=\{(a,b)\in\Z^2:b\equiv ma\pmod p\}.
\]
It has basis $(1,m),(0,p)$ and covolume $p$.
The open disc $a^2+b^2<2p$ has area $2\pi p>4p$, so Minkowski's
theorem gives a nonzero lattice point in it. At that point,
$a^2+b^2\equiv a^2(1+m^2)\equiv0\pmod p$.
The positive integer $a^2+b^2$ is less than $2p$ and divisible by $p$;
therefore it equals $p$.
\end{proof}

\begin{figure}[ht]
\centering
\fig{geometry-ch09-twosquares}
\caption{The proof of Theorem~\ref{geometry:thm:twosquares} for $p=13$, $m=5$: the lattice $b\equiv5a\pmod{13}$ and the disk $a^2+b^2<26$. The four rust points are
$\pm(2,-3)$ and $\pm(3,2)$, and $13=2^2+3^2$.}
\label{geometry:fig:twosquares}
\end{figure}

For $p=13$, choose $m=5$, since $25\equiv-1\pmod{13}$.
The lattice contains $(2,-3)$, giving $13=2^2+3^2$.
The proof does not search each coordinate separately: it arranges a
congruence so that any sufficiently short lattice vector must solve the
arithmetic problem. Exercise~\ref{geometry:ex:foursquares} uses this
same idea in four dimensions.

\section{Plane lattices and Gauss reduction}

A basis can be much longer than the lattice's shortest vector.
Subtracting an integer multiple of one basis vector from the other keeps
the lattice but can shorten the basis substantially.

\begin{definition}
A basis $(b_1,b_2)$ of a plane lattice is reduced if
$|b_1|\le|b_2|$ and $|b_1\cdot b_2|\le|b_1|^2/2$.
\end{definition}

The second inequality says that the component of $b_2$ along $b_1$
is at most half of $b_1$ in either direction.

\begin{proposition}\label{geometry:prop:reducedexists}
Every plane lattice admits a reduced basis.
\end{proposition}
\begin{proof}
Choose a shortest nonzero vector $b_1$. In an existing basis $e_1,e_2$,
write $b_1=pe_1+qe_2$. The integers $p,q$ are coprime, since a common
divisor greater than one would give a shorter nonzero lattice vector.
B\'ezout's identity gives $r,s$ with $ps-qr=1$.
Thus $b_1,c$ with $c=re_1+se_2$ form a basis by
Proposition~\ref{geometry:prop:latticebases}.
Choose an integer $k$ nearest to $(c\cdot b_1)/|b_1|^2$ and set
$b_2=c-kb_1$. The basis is unchanged as a lattice, and the rounding
error gives $|b_1\cdot b_2|\le|b_1|^2/2$.
Also $|b_2|\ge|b_1|$ because $b_2$ is a nonzero lattice vector.
\end{proof}

This proof starts with a shortest vector. Gauss's algorithm finds a
reduced basis without knowing that vector in advance: subtract the
nearest multiple of the shorter vector from the longer, swap their roles
if needed, and repeat. The stopping rule and proof of termination are
in Exercise~\ref{geometry:ex:gaussterminates}.

\begin{figure}[ht]
\centering
\fig{geometry-ch09-reduction}
\caption{Two bases of the same lattice: a long thin one (rust) and the reduced one found by Gauss's algorithm (ochre). Both parallelograms have the same
area.}
\label{geometry:fig:reduction}
\end{figure}

\begin{theorem}\label{geometry:thm:hermite2}
If $(b_1,b_2)$ is a reduced basis of a plane lattice $\Lambda$, then $b_1$
is a shortest nonzero vector and
\[
 \lambda_1(\Lambda)^2\le\frac2{\sqrt3}\det\Lambda.
\]
Equality holds exactly when $\Lambda$ is hexagonal, meaning that it has
a basis of equal-length vectors making angle $60^\circ$.
\end{theorem}
\begin{proof}
Let $x=ub_1+vb_2\ne0$ with $u,v\in\Z$.
If $v=0$, then $|x|=|u||b_1|\ge|b_1|$.
If $v\ne0$, the reduced-basis bounds give
\[
 |x|^2\ge|b_1|^2(u^2-|uv|+v^2).
\]
Indeed, bound the mixed term below by $-|uv||b_1|^2$ and use
$|b_2|^2\ge|b_1|^2$.
Now
$u^2-|uv|+v^2=(|u|-|v|/2)^2+3v^2/4>0$.
It is an integer, so it is at least one. Thus $b_1$ is shortest.

The squared area of its fundamental parallelogram is
\[
 (\det\Lambda)^2
 =|b_1|^2|b_2|^2-(b_1\cdot b_2)^2
 \ge |b_1|^4-\tfrac14|b_1|^4=\tfrac34\lambda_1^4.
\]
Equality requires $|b_2|=|b_1|$ and
$|b_1\cdot b_2|=|b_1|^2/2$, so the angle is $60^\circ$ or
$120^\circ$. Reversing one basis vector changes the latter to the
former. Conversely a hexagonal basis of unit lengths has shortest
length one by the argument above and covolume $\sqrt3/2$.
Rescaling preserves the equality, giving precisely the stated cases.
\end{proof}

For the positive-definite form $ax^2+bxy+cy^2$ this says that its value
at some nonzero integer pair is at most $\sqrt{(4ac-b^2)/3}$.
That arithmetic bound is the same area estimate written in Gram coordinates.

\section{Packings}

Balls of radius $r$ centred at lattice points have disjoint interiors
exactly when $2r\le\lambda_1$. At the largest possible radius
$r=\lambda_1/2$, the density is
\[
 \Delta(\Lambda)=\frac{V_n\lambda_1^n}{2^n\det\Lambda},
\]
where $V_n$ is the unit-ball volume. One fundamental cell accounts for
one ball's volume, even if the ball crosses several cell boundaries.
For example, $V_2=\pi$, $V_3=4\pi/3$, $V_4=\pi^2/2$ and
$V_8=\pi^4/24$.

\begin{corollary}\label{geometry:cor:hexdensity}
Every plane lattice packing of equal discs has density at most
$\pi/\sqrt{12}$. Equality is attained exactly by hexagonal packings.
\end{corollary}
\begin{proof}
Even at the largest radius its density is
\[
 \frac{\pi\lambda_1^2}{4\det\Lambda}
 \le\frac\pi4\frac2{\sqrt3}=\frac\pi{\sqrt{12}}.
\]
The equality cases in Theorem~\ref{geometry:thm:hermite2} give exactly
the hexagonal lattice, with touching neighbouring discs.
\end{proof}

\begin{remark}
Thue's planar packing theorem also bounds nonlattice packings by this
density; its full proof is not given here. In three dimensions,
Gauss's lattice theorem and Hales's unrestricted theorem give the density
$\pi/\sqrt{18}$. We next construct a lattice attaining it, without
claiming that its construction proves either optimality theorem.
\end{remark}

\section{Root lattices and \texorpdfstring{$E_8$}{E8}}

For $n\ge3$, define
\[
 D_n=\{x\in\Z^n:x_1+\cdots+x_n\text{ is even}\}.
\]
The coordinate-sum map to $\Z/2$ is onto, so this subgroup has index two.
The finite-index covolume formula, proved in
Exercise~\ref{geometry:ex:index}, gives $\det D_n=2$.
For clarity, the needed formula has an elementary volume proof:
if a lattice $\Gamma'$ is the union of $k$ cosets of a sublattice
$\Gamma$, count points in a large cube by translating fundamental cells.
The cells crossing the cube's boundary occupy only $O(R^{n-1})$ volume,
so the leading counts give $\det\Gamma'=\det\Gamma/k$.
For existence of a lattice basis, Lagrange's theorem gives
$k\Gamma'\subset\Gamma$, so $\Gamma'$ is a subgroup of the lattice
$\Gamma/k$. In lattice coordinates the subgroup theorem for $\Z^n$
(Volume~3, Lemma~7.5) makes it free of rank $n$, since it contains a
full lattice. A group basis spans the real space and therefore is a real
basis as well. Thus the formula is available for the constructions here.

A nonzero vector in $D_n$ cannot have squared norm one: such a vector
would be $\pm e_i$ and have odd coordinate sum.
Squared norm two is attained precisely by $\pm e_i\pm e_j$, $i<j$.
There are $4\binom n2=2n(n-1)$ such vectors. Hence
\[
 \Delta(D_n)=\frac{V_n}{2^{n/2+1}}.
\]
For $D_3$ this is $\pi/\sqrt{18}\approx0.7405$, and each ball touches
twelve neighbours: this is the face-centred cubic packing.
For $D_4$ the density is $\pi^2/16\approx0.617$, with twenty-four
neighbours.

The half-integer point $h=(1/2,\ldots,1/2)$ is at distance
$\sqrt n/2$ from $D_n$. Each coordinate contributes at least $1/4$
to the squared distance, and the zero vector attains the bound.
For $n=8$ the distance is $\sqrt2$, just the minimum separation in
$D_8$. A translated copy can therefore fill these holes without reducing
that separation.

\begin{theorem}\label{geometry:thm:e8}
Let $h=(1/2,\ldots,1/2)\in\R^8$ and $E_8=D_8\cup(D_8+h)$.
Then:
\begin{enumerate}
\item $E_8$ is a lattice of covolume one;
\item $x\cdot y\in\Z$ and $x\cdot x\in2\Z$ for all $x,y\in E_8$;
\item its shortest nonzero vectors have squared norm two, and there are
exactly $240$ of them.
\end{enumerate}
\end{theorem}
\begin{proof}
(1) The vector $2h=(1,\ldots,1)$ lies in $D_8$.
The two cosets therefore form a group: adding two elements of $D_8+h$
lands in $D_8$, and negating that coset gives the same coset.
It is a finite-index overgroup of $D_8$, of index two, so the formula
above makes it a lattice of covolume $2/2=1$.

(2) If $x\in D_8$, then
$|x|^2=\sum x_i^2\equiv\sum x_i\equiv0\pmod2$.
If $x=n+h$ with $n\in D_8$, then
\[
 |x|^2=\sum_i(n_i^2+n_i+1/4)=\sum_i n_i(n_i+1)+2\in2\Z.
\]
Since $x+y\in E_8$, polarisation gives
$2x\cdot y=|x+y|^2-|x|^2-|y|^2\in2\Z$.
Thus $x\cdot y$ is integral.

(3) In $D_8$, the squared-norm-two vectors are the $112$ vectors
$\pm e_i\pm e_j$. In $D_8+h$, every coordinate has absolute value
at least $1/2$, giving squared norm at least $8/4=2$.
Equality requires all coordinates to be $\pm1/2$.
Writing them as $n_i+1/2$ makes $n_i$ either $0$ or $-1$;
the condition $n\in D_8$ requires an even number of minus signs.
Exactly half of the $2^8$ sign choices qualify, giving $128$ vectors.
The total is $112+128=240$.
\end{proof}

A lattice whose dot products are integral is called integral; it is even
if every squared norm is even. Covolume one is called unimodular.
Thus $E_8$ is even and unimodular, with density
\[
 \Delta(E_8)=\frac{\pi^4}{24}\frac{(\sqrt2)^8}{2^8}
 =\frac{\pi^4}{384}\approx0.2537.
\]
This is twice the density of $D_8$, and each ball touches $240$ neighbours.

\begin{figure}[ht]
\centering
\fig{geometry-ch09-e8}
\caption{The $240$ shortest vectors of $E_8$, projected onto a well-chosen plane (the \emph{Coxeter plane}): they fall on eight circles of $30$ points each.}
\label{geometry:fig:e8}
\end{figure}

\begin{remark}
Viazovska's optimality proof constructs a function from modular forms
whose Fourier transform helps bound every packing's density.
Cohn, Kumar, Miller, Radchenko and Viazovska obtained the corresponding
optimality theorem in dimension twenty-four for the Leech lattice.
Those theorems are quoted, not proved by the root counts above.
The restriction that positive-definite even unimodular lattices have
dimension divisible by eight is also a quoted theorem; Serre's book in
the sources proves it with modular forms. The same $240$ vectors form
the roots associated with the exceptional Lie group of type $E_8$.
\end{remark}

\begin{figure}[ht]
\centering
\fig{v9-app-h-folding}
\caption{Blichfeldt's overlap step for $\Z^2$. Two small pieces lie in adjacent unit cells. Translating the rust piece by $-(1,0)$ into the half-open fundamental square makes it overlap the charcoal piece. A set of area greater than one forces an overlap somewhere, even when its pieces have different shapes.}
\label{geometry:fig:v9-folding}
\end{figure}
\begin{figure}[ht]
\centering
\fig{v9-app-h-basis-reduction}
\caption{Subtracting twice $u=(1,0)$ from $v=(2.4,1)$ gives $v'=(0.4,1)$. The rust and charcoal parallelograms have the same area, one, and generate the same lattice. The new basis is reduced: $|u|\le|v'|$ and $|u\cdot v'|=0.4\le0.5$.}
\label{geometry:fig:v9-basis-reduction}
\end{figure}

A packing picture hides two separate tasks: constructing a dense arrangement
and proving that no arrangement is denser. Determinants and reduction settle
both tasks for plane lattices. In higher dimensions, even a short explicit
construction such as $E_8$ can need a very different argument for optimality.
Appendix~I returns to distance itself and compares the sphere with the
hyperbolic plane.

\section{Exercises}

\begin{exercise}[$\star$]\label{geometry:ex:basisZ2}
Decide whether each pair $(1,2),(3,4)$ and $(2,3),(3,5)$ is a basis of
$\Z^2$. Describe the generated lattice in each case.
\end{exercise}
\begin{exercise}[$\star$]\label{geometry:ex:index}
Suppose $\Lambda\subset\Lambda'\subset\R^n$, where $\Lambda$ is a
lattice and $\Lambda'$ is a subgroup with $[\Lambda':\Lambda]=k<\infty$.
Prove that $\Lambda'$ is a lattice with covolume $\det\Lambda/k$.
Apply the result to $D_n\subset\Z^n$.
\end{exercise}
\begin{exercise}[$\star$]\label{geometry:ex:gaussreduce}
Apply Gauss reduction to the basis $(13,8),(21,13)$ of $\Z^2$.
Identify the number sequence in the vectors and explain the relation to
the ordinary Euclidean algorithm.
\end{exercise}
\begin{exercise}[$\star\star$]\label{geometry:ex:gaussterminates}
Give a precise stopping rule for Gauss's algorithm and prove that it
terminates with a reduced basis.
\end{exercise}
\begin{exercise}[$\star\star$]\label{geometry:ex:foursquares}
Let $p$ be an odd prime. First find integers $a,b$ with
$a^2+b^2+1\equiv0\pmod p$. For the lattice
\[
 \{x\in\Z^4:x_3\equiv ax_1+bx_2,\quad
 x_4\equiv bx_1-ax_2\pmod p\},
\]
find a nonzero vector with squared norm less than $2p$ and deduce that
$p$ is a sum of four squares. You may use the four-ball volume
$\pi^2r^4/2$.
\end{exercise}
\begin{exercise}[$\star$]\label{geometry:ex:kissing2}
Prove that a disc can touch at most six equal-radius discs whose interiors
are pairwise disjoint and disjoint from its own interior.
\end{exercise}
\begin{exercise}[$\star\star$]\label{geometry:ex:fcc}
Show that $(1,1,0),(1,0,1),(0,1,1)$ generate $D_3$ and count its nearest
neighbours. Describe the layers $x_3=0$ and $x_3=1$ as shifted square grids.
\end{exercise}
\begin{exercise}[$\star\star$]\label{geometry:ex:deephole}
For $n\ge3$, find the distance from $h=(1/2,\ldots,1/2)$ to $D_n$.
For even $n$, prove $D_n\cup(D_n+h)$ is a lattice and that its minimum
nonzero squared norm is $\min(2,n/4)$. For which such $n$ does its
packing density exceed that of $D_n$?
\end{exercise}
\begin{exercise}[$\star\star$]\label{geometry:ex:minkowskibound}
Use Minkowski's theorem with a disc to prove
$\lambda_1^2\le(4/\pi)\det\Lambda$ for a plane lattice.
Compare the constant with the reduced-basis bound.
\end{exercise}
\begin{exercise}[$\star\star$]\label{geometry:ex:circlecount}
Let $N(r)=\#(\Z^2\cap\overline D(0,r))$.
For $r\ge1$, prove
$|N(r)-\pi r^2|\le2\sqrt2\pi r+2\pi$ and deduce $N(r)/r^2\to\pi$.
\end{exercise}
\begin{exercise}[$\star$]\label{geometry:ex:e8roots}
Show that the $240$ shortest vectors generate $E_8$ as an additive group.
If two such vectors satisfy $x\ne\pm y$, prove
$x\cdot y\in\{-1,0,1\}$.
\end{exercise}
\begin{exercise}[$\star\star\star$]\label{geometry:ex:hermite}
Prove by induction on $n\ge1$ that every lattice satisfies
\[
 \lambda_1^2\le(4/3)^{(n-1)/2}(\det\Lambda)^{2/n}.
\]
\end{exercise}

\begin{exercise}[$\star\star$]\label{geometry:ex:v9-closed-minkowski}
\emph{The boundary case.} Let $C\subset\R^n$ be compact, convex and
symmetric, with volume at least $2^n\det\Lambda$.
Show that $C$ contains a nonzero lattice point. Explain precisely why the
same conclusion fails for the open cube at the equality threshold.
\end{exercise}
\begin{exercise}[$\star\star$]\label{geometry:ex:v9-dual-lattice}
\emph{Integer measurements of a lattice.} Define
$\Lambda^*=\{y\in\R^n:y\cdot x\in\Z\text{ for every }x\in\Lambda\}$.
Show that $(B\Z^n)^*=B^{-\mathsf T}\Z^n$, that the two covolumes
multiply to one, and that $(\Lambda^*)^*=\Lambda$.
Use integrality and unimodularity to prove $E_8^*=E_8$.
\end{exercise}
\begin{exercise}[$\star\star$]\label{geometry:ex:v9-complex-lattice}
\emph{A sublattice from one complex integer.} Identify $\Z^2$ with
$\Z[i]$. For nonzero $\alpha=a+bi\in\Z[i]$, find a real basis,
covolume, index in $\Z[i]$ and shortest length for $\alpha\Z[i]$.
Determine its largest-disc packing density. Explain why a large index
does not improve that density.
\end{exercise}
\begin{exercise}[$\star\star$]\label{geometry:ex:v9-reduction-domain}
\emph{A lattice in modular coordinates.} Identify an oriented plane with
$\C$ and choose a positively oriented basis $b_1,b_2$ of a lattice.
Put $\tau=b_2/b_1\in\mathbb H$.
Show that reduction is equivalent to $|\operatorname{Re}\tau|\le1/2$
and $|\tau|\ge1$. Describe the changes in $\tau$ when $b_2$ is replaced
by $b_2-kb_1$, or the basis by $(b_2,-b_1)$.
Recover the sharp Hermite bound from
$\det\Lambda=|b_1|^2\operatorname{Im}\tau$ and classify equality.
\end{exercise}

\section{Hints}

\begin{hint}{geometry:ex:basisZ2}
Compute each determinant and compare with $\det\Z^2=1$.
\end{hint}
\begin{hint}{geometry:ex:index}
Lagrange's theorem gives $k\Lambda'\subset\Lambda$.
Then use the subgroup theorem for $\Z^n$ and count cosets in large cubes.
\end{hint}
\begin{hint}{geometry:ex:gaussreduce}
Round the first ratio $377/233$ to $2$; the next is $-89/34$.
\end{hint}
\begin{hint}{geometry:ex:gaussterminates}
Stop when the projection coefficient has absolute value at most $1/2$.
Otherwise replacing the longer vector strictly decreases its length.
\end{hint}
\begin{hint}{geometry:ex:foursquares}
Two sets of $(p+1)/2$ quadratic-residue values must overlap.
The two congruences define a lattice of index $p^2$.
\end{hint}
\begin{hint}{geometry:ex:kissing2}
Neighbouring centres on the radius-two circle must have angular separation
at least $\pi/3$.
\end{hint}
\begin{hint}{geometry:ex:fcc}
Solve for the three generator coefficients in terms of an even-sum vector.
\end{hint}
\begin{hint}{geometry:ex:deephole}
Each half-integer coordinate contributes at least $1/4$ to a squared norm.
Compare $2(\min(2,n/4)/2)^{n/2}$ with one.
\end{hint}
\begin{hint}{geometry:ex:minkowskibound}
The open disc of radius $\lambda_1$ has no nonzero lattice point.
\end{hint}
\begin{hint}{geometry:ex:circlecount}
Attach half-open unit squares to the counted points. Use inner radius
$\max(0,r-\sqrt2)$, not a negative radius.
\end{hint}
\begin{hint}{geometry:ex:e8roots}
Generate $D_8$ by its norm-two vectors, then adjoin $h$.
For dot products, combine integrality with strict Cauchy--Schwarz.
\end{hint}
\begin{hint}{geometry:ex:hermite}
Extend a shortest vector to a lattice basis, project to its perpendicular
space, and round a lift along that shortest direction.
\end{hint}
\begin{hint}{geometry:ex:v9-closed-minkowski}
Apply the strict theorem to $(1+1/j)C$ and use finiteness of lattice
points in a bounded set.
\end{hint}
\begin{hint}{geometry:ex:v9-dual-lattice}
The defining condition is $B^{\mathsf T}y\in\Z^n$.
Use the finite-index covolume formula for $E_8\subset E_8^*$.
\end{hint}
\begin{hint}{geometry:ex:v9-complex-lattice}
Multiplication by $a+bi$ has real matrix
$\left(\begin{smallmatrix}a&-b\\b&a\end{smallmatrix}\right)$.
\end{hint}
\begin{hint}{geometry:ex:v9-reduction-domain}
The two ratios are $|b_2|/|b_1|=|\tau|$ and
$(b_1\cdot b_2)/|b_1|^2=\operatorname{Re}\tau$.
\end{hint}

\section{Solutions}

\begin{solution}{geometry:ex:basisZ2}
\idea The determinant measures how many integer points a fundamental cell
accounts for.
\writeup The first pair has determinant $-2$, so does not generate
$\Z^2$. Its vectors are $(p+3q,2p+4q)$, whose second coordinate is
even. Conversely, $(1,2)$ and $(0,-2)=(3,4)-3(1,2)$ generate every
$(a,2b)$, so this is exactly the sublattice with even second coordinate.
The second pair has determinant $2\cdot5-3\cdot3=1$.
Its integer inverse makes it a basis of $\Z^2$.
\end{solution}
\begin{solution}{geometry:ex:index}
\idea A finite-index overgroup sits inside a fixed rescaling of the original
lattice, so it remains discrete.
\writeup For $x\in\Lambda'$, the class of $x$ in the order-$k$ group
$\Lambda'/\Lambda$ satisfies $kx\in\Lambda$.
Thus $\Lambda\subset\Lambda'\subset\Lambda/k$.
In coordinates $kB^{-1}$ this is a subgroup of $\Z^n$ containing
$k\Z^n$. The subgroup theorem (Volume~3, Lemma~7.5) makes it free;
it has rank $n$, since it contains $n$ real-independent vectors and
cannot have greater rank as a subgroup of $\Z^n$.
A group basis spans $\R^n$, because the subgroup contains $k\Z^n$,
so it is also real-independent and gives a lattice basis.

For any lattice $\Gamma$, bounded fundamental cells imply
\[
 \#(\Gamma\cap[-R,R]^n)
 =\frac{(2R)^n}{\det\Gamma}+O(R^{n-1}).
\]
Indeed, cells based well inside the cube are contained in it, and only
cells in a fixed-width boundary layer cause a discrepancy.
The same estimate holds for a fixed translated coset.
Since $\Lambda'$ is a union of $k$ cosets of $\Lambda$, compare the
leading terms to obtain $1/\det\Lambda'=k/\det\Lambda$.
The coordinate-sum homomorphism gives $[\Z^n:D_n]=2$, so
$\det D_n=2\det\Z^n=2$.
\end{solution}
\begin{solution}{geometry:ex:gaussreduce}
\idea Nearly parallel Fibonacci vectors reduce to much shorter Fibonacci
vectors.
\writeup Begin with $u=(13,8)$ and $v=(21,13)$.
The projection coefficient is $377/233$, whose nearest integer is $2$.
Thus $v-2u=(-5,-3)$. After swapping, the coefficient is
$-89/34$, which rounds to $-3$; the next vector is
$(13,8)+3(-5,-3)=(-2,-1)$.
Next $(-5,-3)-3(-2,-1)=(1,0)$, and then
$(-2,-1)+2(1,0)=(0,-1)$.
The final basis $(1,0),(0,-1)$ is reduced.
The coordinates $21,13,8,5,3,2,1$ are consecutive Fibonacci numbers.
The long initial vectors have nearly equal slopes, so their differences
are much shorter. Ordinary Euclidean division follows the same sequence
with single subtractions; nearest-integer rounding combines some of those
subtractions. The first two projection ratios are close to the golden
ratio and minus its square, respectively.
\end{solution}
\begin{solution}{geometry:ex:gaussterminates}
\idea Every nonterminal round moves within a finite set while strictly
lowering a length sum.
\writeup First order the basis so $|u|\le|v|$, and put
$\mu=(u\cdot v)/|u|^2$. If $|\mu|\le1/2$, stop.
Otherwise choose an integer $k$ nearest to $\mu$, replace $v$ by
$v-ku$, and reorder. Write $v=\mu u+w$ with $w\perp u$.
Then
$|v-ku|^2=(\mu-k)^2|u|^2+|w|^2<\mu^2|u|^2+|w|^2=|v|^2$,
since $|\mu-k|\le1/2<|\mu|$.
All replacements preserve the lattice basis and independence.
The sum of squared lengths strictly decreases at each round, and all
vectors remain no longer than the initial longer vector.
There are finitely many lattice vectors in that ball and hence finitely
many possible pairs. The process must stop, precisely at a reduced basis.
The stopping rule avoids cycling when $|\mu|=1/2$.
\end{solution}
\begin{solution}{geometry:ex:foursquares}
\idea Congruences make the norm a multiple of $p$; a volume bound makes
that multiple smaller than two.
\writeup The sets $\{a^2\}$ and $\{-1-b^2\}$ each have
$(p+1)/2$ residues and must overlap, giving the desired $a,b$.
The map from $\Z^4$ to $(\Z/p)^2$ defined by the two congruence errors
is onto by varying $x_3,x_4$. Its kernel has index $p^2$ and covolume
$p^2$; an explicit basis is
$(1,0,a,b),(0,1,b,-a),(0,0,p,0),(0,0,0,p)$.
On that kernel,
\[
 x_3^2+x_4^2\equiv(a^2+b^2)(x_1^2+x_2^2)
 \equiv-(x_1^2+x_2^2)\pmod p.
\]
The open ball $|x|^2<2p$ has volume $2\pi^2p^2>16p^2$.
Minkowski gives a nonzero kernel vector in it, with
$0<|x|^2<2p$ and $p\mid|x|^2$, hence $|x|^2=p$.
Together with $2=1^2+1^2+0^2+0^2$ and Euler's four-square identity,
this also gives the usual extension from primes to all positive integers.
The identity itself was established earlier in the series' quaternion
arithmetic; the prime assertion just proved does not depend on that extension.
\end{solution}
\begin{solution}{geometry:ex:kissing2}
\idea Each neighbouring centre occupies at least a sixty-degree angular
interval around the central disc.
\writeup Scale the radius to one and put the central centre at $0$.
Touching centres lie on $|x|=2$. Distinct centres have separation at
least two. If their smaller angular separation is $\theta$, then
$|c_i-c_j|=4\sin(\theta/2)\ge2$, so $\theta\ge\pi/3$.
The cyclic gaps between consecutive centres therefore each have angle
at least $\pi/3$ (a gap larger than $\pi$ also meets this bound).
Their sum is $2\pi$, allowing at most six neighbours.
The regular hexagonal arrangement attains six.
\end{solution}
\begin{solution}{geometry:ex:fcc}
\idea Even coordinate sum makes the generator coefficients integral.
\writeup For $x=(x_1,x_2,x_3)\in D_3$, set
\[
 A=(x_1+x_2-x_3)/2,\quad B=(x_1+x_3-x_2)/2,\quad
 C=(x_2+x_3-x_1)/2.
\]
All three are integers and
$x=A(1,1,0)+B(1,0,1)+C(0,1,1)$.
The generators are independent, with determinant $-2$, so form a basis.
The twelve nearest neighbours are $\pm e_i\pm e_j$.
At height zero, $x_1+x_2$ is even, giving a square grid of spacing
$\sqrt2$ rotated by $45^\circ$. At height one the sum is odd, giving
that same grid shifted by $(1,0)$ to the square centres of the lower
layer. This is the square-layer view of the close packing.
\end{solution}
\begin{solution}{geometry:ex:deephole}
\idea Coordinatewise distance gives a bound attained by an allowed parity
choice.
\writeup For every integer vector $x$, each $|x_i-1/2|\ge1/2$,
so $|x-h|^2\ge n/4$. The vector $0\in D_n$ attains this bound.
For even $n$, $2h\in D_n$, so the two cosets form a group of covolume
one, as in the $E_8$ proof. In its integer coset the minimum squared
norm is two. In its half-integer coset it is $n/4$, attained by $h$.
Thus the minimum is $\min(2,n/4)$.
The density ratio to $D_n$ is
$2(\min(2,n/4)/2)^{n/2}$.
For even $n\ge8$ it equals two. For $n=4$ it is $1/2$, and for
$n=6$ it is $2(3/4)^3=27/32<1$.
These are all even $n\ge3$, so improvement occurs exactly for even
$n\ge8$.
\end{solution}
\begin{solution}{geometry:ex:minkowskibound}
\idea Apply the contrapositive to the open disc just short of admitting a
shortest vector.
\writeup The open disc of radius $\lambda_1$ contains no nonzero
lattice point. Minkowski therefore forces its area to be at most
$4\det\Lambda$. Hence
$\lambda_1^2\le(4/\pi)\det\Lambda$.
The constant $4/\pi\approx1.273$ exceeds $2/\sqrt3\approx1.155$,
so reduction gives the sharper planar bound.
\end{solution}
\begin{solution}{geometry:ex:circlecount}
\idea Unit squares compare a discrete count with nearby circular areas.
\writeup Attach the half-open square $[a,a+1)\times[b,b+1)$ to every
counted point $(a,b)$. Their union has area $N(r)$ and lies inside the
disc of radius $r+\sqrt2$. If $r\ge\sqrt2$, it contains the disc of
radius $r-\sqrt2$: the lower-left integer corner of the square containing
a point of that disc is at distance at most $r$ from $0$.
Thus for $r\ge\sqrt2$,
\[
 \pi(r-\sqrt2)^2\le N(r)\le\pi(r+\sqrt2)^2,
\]
which implies the desired error bound.
For $1\le r<\sqrt2$, the same upper bound holds and $N(r)\ge0$.
The possible lower error $\pi r^2-N(r)$ is at most $\pi r^2$,
which is smaller than $2\sqrt2\pi r+2\pi$ in this range.
This completes the stated range without assigning a negative inner radius.
After division by $r^2$, the error tends to zero, proving the limit.
The sharper order of the error is Gauss's circle problem, not resolved by
this boundary-strip estimate.
\end{solution}
\begin{solution}{geometry:ex:e8roots}
\idea Integer dot products and a strict length bound leave only three
possibilities.
\writeup For shortest vectors $x\ne\pm y$, strict Cauchy--Schwarz gives
$|x\cdot y|<|x||y|=2$. Integrality therefore leaves $-1,0,1$.
To prove generation, pair the odd coordinates of any vector in $D_8$;
their number is even. Subtract vectors $\pm e_i\pm e_j$ to make all
coordinates even. Each remaining $2e_i$ is
$(e_i+e_j)+(e_i-e_j)$, again a sum of roots.
Thus the integer roots generate $D_8$. The vector $h$ is also a shortest
vector and adjoining it generates $D_8+\Z h=E_8$.
\end{solution}
\begin{solution}{geometry:ex:hermite}
\idea Project away a shortest direction, then lift a short projected vector
with a small component in the removed direction.
\writeup In dimension one, the formula is equality.
For $n\ge2$, choose a shortest vector $b_1$.
Every vector of $\Lambda\cap\R b_1$ is an integer multiple of $b_1$:
otherwise subtracting the nearest lower multiple leaves a nonzero
shorter vector on that line. The quotient $\Lambda/\Z b_1$ is
torsion-free. Indeed, if $kx=mb_1$ with $k\ne0$, then $x$ lies on
$\R b_1$, so is already an integer multiple of $b_1$.
It is a finitely generated abelian group, so Volume~3, Theorem~7.6,
makes it free of rank $n-1$. Lift a quotient basis to $b_2,\ldots,b_n$;
together these vectors form a group basis of $\Lambda$.

Project orthogonally to $b_1^\perp$ by $\pi$.
The vectors $\pi(b_2),\ldots,\pi(b_n)$ are real-independent:
a real relation among them would make a real combination of
$b_2,\ldots,b_n$ a multiple of $b_1$, contradicting the real basis
property. They generate $\pi(\Lambda)$, so it is a lattice.
The height-times-base determinant formula gives
$\det\pi(\Lambda)=\det\Lambda/|b_1|$.
By induction there is a nonzero $y=\pi(x)$ with
\[
 |y|^2\le(4/3)^{(n-2)/2}
 \left(\frac{\det\Lambda}{|b_1|}\right)^{2/(n-1)}.
\]
Change $x$ by a multiple of $b_1$ so its parallel component has length
at most $|b_1|/2$. The vector is nonzero because $y\ne0$, so
$|b_1|^2\le|x|^2\le|b_1|^2/4+|y|^2$.
Hence $3|b_1|^2/4\le|y|^2$. Combining and rearranging gives
\[
 |b_1|^{2n/(n-1)}\le(4/3)^{n/2}(\det\Lambda)^{2/(n-1)}.
\]
Raise to the power $(n-1)/n$ to obtain the result.
\end{solution}

\begin{solution}{geometry:ex:v9-closed-minkowski}
\idea Enlarge slightly to enter the strict theorem, then keep a repeated
lattice point as the enlargement disappears.
\writeup The positive threshold ensures $C$ is nonempty and has positive
volume. Symmetry and convexity imply $0\in C$, so
$(1+1/j)C\subset2C$ for $j\ge1$.
Each enlarged set has volume strictly greater than $2^n\det\Lambda$,
so contains a nonzero lattice vector $v_j$.
The bounded set $2C$ contains finitely many lattice points; some nonzero
$v$ occurs for infinitely many indices $j_k\to\infty$.
Then $v/(1+1/j_k)\in C$ and tends to $v$.
Closedness gives $v\in C$.
For $C=(-1,1)^n$ and $\Lambda=\Z^n$, the volume is exactly $2^n$
but all nonzero integer points are excluded. The limiting vector can lie
on a missing boundary, so that last closedness step fails.
\end{solution}
\begin{solution}{geometry:ex:v9-dual-lattice}
\idea Test against the basis vectors rather than every lattice point.
\writeup For $x=Bm$, the condition $y\cdot Bm\in\Z$ for all
$m\in\Z^n$ is equivalent to $B^{\mathsf T}y\in\Z^n$.
Thus $\Lambda^*=B^{-\mathsf T}\Z^n$, a lattice of covolume
$|\det B|^{-1}$. Applying the formula twice gives $\Lambda^{**}=\Lambda$.
Integrality of $E_8$ means $E_8\subset E_8^*$.
Both have covolume one. The inclusion has finite index: expressing a
basis of $E_8$ in a basis of $E_8^*$ gives a nonsingular integer matrix,
whose adjugate shows the quotient is finite. The covolume formula makes
that index one. Hence $E_8=E_8^*$.
\end{solution}
\begin{solution}{geometry:ex:v9-complex-lattice}
\idea Complex multiplication is a rotation followed by a uniform dilation.
\writeup Multiplying the basis $1,i$ by $\alpha$ gives
$(a,b),(-b,a)$, with determinant $a^2+b^2>0$.
Therefore the covolume and index are both $a^2+b^2$.
Every nonzero Gaussian integer has modulus at least one, attained by
$1,-1,i,-i$. Since $|\alpha z|=|\alpha||z|$, the shortest length is
$|\alpha|=\sqrt{a^2+b^2}$, with four shortest vectors.
The largest packing radius is $|\alpha|/2$, giving density
\[
 \frac{\pi|\alpha|^2}{4(a^2+b^2)}=\frac\pi4.
\]
Increasing the index dilates the cell and the discs by matching area
factors. It changes their scale, not the square shape of the packing.
\end{solution}
\begin{solution}{geometry:ex:v9-reduction-domain}
\idea Divide by the first vector to turn the lattice basis conditions into
conditions on one complex number.
\writeup Complex division by $b_1$ is a rotation and dilation.
Thus $|b_2|=|b_1||\tau|$ and
$b_1\cdot b_2=|b_1|^2\operatorname{Re}\tau$.
The two reduced conditions are exactly $|\tau|\ge1$ and
$|\operatorname{Re}\tau|\le1/2$.
Subtracting $kb_1$ gives $\tau\mapsto\tau-k$.
The positively oriented basis $(b_2,-b_1)$ gives
$\tau\mapsto-1/\tau$; both operations preserve the upper half-plane.

The area formula is
$\det\Lambda=|b_1|^2\operatorname{Im}\tau$.
For reduced $\tau$, its imaginary part is at least $\sqrt3/2$, since
$|\tau|\ge1$ and $|\operatorname{Re}\tau|\le1/2$.
The theorem above identifies $|b_1|=\lambda_1$, giving
$\lambda_1^2\le2\det\Lambda/\sqrt3$.
Equality forces $\tau=1/2+i\sqrt3/2$ or $-1/2+i\sqrt3/2$.
These describe equal-length bases of angle $60^\circ$ or $120^\circ$,
hence exactly the hexagonal lattice. This is the modular-domain picture
of the same sharp bound, not a separate numerical approximation.
\end{solution}

\section*{Sources and further reading}

\begin{itemize}
\item T.~C.~Hales, \emph{A computer verification of the Kepler conjecture}, arXiv:math/0305012 (2003) (Harriot, Raleigh and Kepler's 1611 essay, quoted in the
epigraph); T.~C.~Hales, \emph{A proof of the Kepler conjecture}, Annals of Mathematics 162 (2005), 1065--1185; T.~C.~Hales et al., \emph{A formal proof of the
Kepler conjecture}, Forum of Mathematics, Pi 5 (2017).
\item J.~Kepler, \emph{The Six-Cornered Snowflake} (1611), English translation, Clarendon Press, Oxford, 1966.
\item M.~Viazovska, \emph{The sphere packing problem in dimension 8}, Annals of Mathematics 185 (2017); H.~Cohn, A.~Kumar, S.~D.~Miller, D.~Radchenko and
M.~Viazovska, \emph{The sphere packing problem in dimension 24}, Annals of Mathematics 185 (2017); D.~de~Laat and F.~Vallentin, \emph{A breakthrough in sphere
packing: the search for magic functions}, arXiv:1607.02111 (an account for a general audience).
\item J.~H.~Conway and N.~J.~A.~Sloane, \emph{Sphere Packings, Lattices and Groups}, Springer, 3rd ed., 1999 (everything about $D_n$, $E_8$ and the Leech lattice).
\item H.~Minkowski, \emph{Geometrie der Zahlen}, Teubner, 1896 and 1910; J.~W.~S.~Cassels, \emph{An Introduction to the Geometry of Numbers}, Springer, 1959.
\item J.-P.~Serre, \emph{A Course in Arithmetic}, Springer, 1973 (even unimodular lattices and modular forms).
\item The story of Hilbert's eighteenth problem and of the planar case (Thue) is told in Hales's article above and in C.~Zong, \emph{Sphere Packings},
Springer, 1999.
\end{itemize}
